\scanpage{23}
Schémas abéliens sur un $S$ \uline{régulier}

a) Un schéma préabélien $P$ sur $S$ est projectif sur $S$
(Raisonnement de \unsure{Nakai|Nakaï}-Shimura)

b) Il existe un schéma de translations associé $G$,
et $P$ est un espace \deleted{fibré} principal homogène
\deleted{\unsure{homotopique|homogène}} isotrivial sous $G$.

c) Si $G$ est un schéma abélien sur $S$, alors
$H^1(S,G)\subset H^1(K,G)$. Si $P_K$ est un
fibré principal homogène sous $G_{\unsure{K|k}}$, pour
qu'il \deleted{soit} provienne d'un $P$ pr. hom
ss $G$, il f. et s. qu'il ``se réduise bien''

\margin{\unsure{?}} En tous cas, il y a un plus grand
\deleted{\illegible{1}} ouvert $U$ \deleted{tel} dense de $S$ tel que
$P_K$ provienne d'un $P$ sur $U$ (c'est
l'ens des pts où $P_K$ \deleted{est} se réduit bien)

% slanted
\margin{résulte du th. de pureté}
d) $S-U=Z$ est de codim $1$ en \unsure{tout|tt} point.

Donc
\[
H^1(K/G)/H^1(S,G)\subset
\sum_{\substack{D\ \text{div}\\\deleted{\text{\unsure{irréductibles|irréductible}}}\\\text{premiers}\\\text{de }S}}
H^1(K_{\unsure{D|v}},G)/H^1(\uline{\unsure{O|\mathcal{O}}}_D,\unsure{G}
\]
\margin{\uline{\unsure{gpe|groupe}}: \uline{\unsure{``étudié''|étudié|étudier}} ne \uline{\unsure{ch.ge|change}} \uline{pas par} \uline{complétion}}

Analogue \inserted{th. de pureté} pour la bonne
réduction des schémas abéliens?
% archive sheet number: 22
